\originalsection{第4次习题课}
\sourceinfo{MIT 18.04, Spring 2018, Recitation 4; 题面 \nolinkurl{mit18_04s18_recit4-handout.pdf}, PDF第1页; 解答 \nolinkurl{mit18_04s18_recit4-solutions.pdf}, PDF第1--2页. 课程作者 Jeremy Orloff, 页内署名 Vishesh Jain, 2018, CC BY-NC-SA 4.0}
\begin{exercise}
\textbf{原题 Rec4.1。}\label{ass:rec4-1}
用 Cauchy 积分公式计算 $I=\int_{-\infty}^{\infty}(1+x^2)^{-2}\dd x$. 可用积分三角不等式 $|\int_\Gamma f(z)\dd z|\le\int_\Gamma|f(z)|\,|\dd z|$. 1.1. 考虑以0为圆心、半径 $R$ 的上半圆闭轮廓 $C$, 计算 $\int_C(1+z^2)^{-2}\dd z$. 1.2. 将 $C=C_1\cup C_2$ 分解为实轴上 $-R$ 到 $R$ 的线段和上半圆弧, 给出圆弧积分绝对值的上界. 1.3. 估计线段积分, 并令 $R\to\infty$.
\end{exercise}
\begin{solution}
编者补充 (AI): 官方解答PDF第1页仅指向讲义例4.11. 取 $R>1$ 且闭轮廓逆时针, 则仅有 $\ii$ 位于内部. 令 $g(z)=(z+\ii)^{-2}$, 由导数型 Cauchy 公式,
\[
\int_C\frac{\dd z}{(1+z^2)^2}=2\pi\ii g'(\ii)
=2\pi\ii\frac{-2}{(2\ii)^3}=\frac\pi2.
\]
在 $C_2$ 上 $|1+z^2|\ge R^2-1$, 弧长为 $\pi R$, 所以圆弧积分的模不超过 $\pi R/(R^2-1)^2\to0$. 于是 $|\int_{-R}^{R}(1+x^2)^{-2}\dd x-\pi/2|\le\pi R/(R^2-1)^2$. 被积函数在有限区间连续且在无穷远为 $O(x^{-4})$, 两个单侧广义积分都收敛; 因而上述对称极限就是通常的广义积分, 得 $I=\pi/2$.
\end{solution}
\begin{exercise}
\textbf{原题 Rec4.2。}\label{ass:rec4-2}
2.1 (Cauchy不等式). 设 $C_R$ 为以 $z_0$ 为圆心、半径 $R$ 的圆, $f$ 在圆及其内部解析. 令 $M_R=\max_{z\in C_R}|f(z)|$, 用导数型积分公式证明 $|f^{(n)}(z_0)|\le n!M_R/R^n$. 2.2 (Liouville定理). 若 $f$ 为整函数且对所有 $z\in\C$ 有 $|f(z)|\le M$, 由 $n=1$ 的情形能推出什么?
\end{exercise}
\begin{solution}
编者补充 (AI): 官方解答PDF第1页仅指向讲义定理4.15与4.16. 圆及其内部解析按通常含义是存在闭圆盘的开邻域, $f$ 在其上解析. 导数型公式与长度估计给出
\[
|f^{(n)}(z_0)|=\left|\frac{n!}{2\pi\ii}\int_{C_R}\frac{f(z)}{(z-z_0)^{n+1}}\dd z\right|
\le\frac{n!}{2\pi}\frac{M_R}{R^{n+1}}(2\pi R)=\frac{n!M_R}{R^n}.
\]
对整函数可任意增大 $R$. 固定 $z_0$ 后 $|f'(z_0)|\le M/R\to0$, 所以 $f'=0$ 处处成立. 沿任一线段应用原函数积分公式得到 $f(b)-f(a)=\int_a^b f'(z)\dd z=0$, 故 $f$ 为常函数.
\end{solution}
\begin{exercise}
\textbf{原题 Rec4.3。}\label{ass:rec4-3}
设 $P(z)=a_nz^n+a_{n-1}z^{n-1}+\cdots+a_0$, $a_n\ne0$, $n\ge1$. 证明 $P$ 在复数域恰有 $n$ 个根, 按重数计. 3.1. 反设 $P$ 没有零点, 证明 $f=1/P$ 整且有界, 用 Liouville 定理导出矛盾. 3.2. 迭代至少一个根的结论, 得到全部根数.
\end{exercise}
\begin{solution}
编者补充 (AI): 官方解答PDF第1页仅指向讲义4.7.2节. 因 $P(z)/(a_nz^n)\to1$, 存在 $R$ 使 $|z|\ge R$ 时 $|P(z)|\ge|a_n||z|^n/2$, 因而 $|1/P(z)|\le2/(|a_n|R^n)$. 闭圆盘 $|z|\le R$ 上, 连续函数 $|P|$ 在无零点假设下具有正最小值 $m$, 所以 $|1/P|\le1/m$. 这证明 $1/P$ 是有界整函数, 从而为常数. 但 $1/P(z)\to0$ 当 $|z|\to\infty$, 所以常数必须为0, 与 $1/P$ 处处非零矛盾. 因此存在根 $\alpha_1$. 多项式除法给 $P=(z-\alpha_1)P_1$, $P_1$ 的次数为 $n-1$. 对非恒定的 $P_1$ 重复, 最终 $P(z)=a_n\prod_{j=1}^n(z-\alpha_j)$. 此乘积分解显示每个零点的重数, 重数和恰为 $n$, 且不存在其他零点.
\end{solution}
\begin{exercise}
\textbf{原题 Rec4.4。}\label{ass:rec4-4}
设 $C_R$ 为以 $z_0$ 为圆心、半径 $R$ 的圆, $f$ 在圆及其内部解析. 用 Cauchy 积分公式证明均值性质 $f(z_0)=(2\pi)^{-1}\int_0^{2\pi}f(z_0+R\ee^{\ii\theta})\dd\theta$.
\end{exercise}
\begin{solution}
编者补充 (AI): 官方解答PDF第2页仅指向讲义定理4.18. 将正向圆参数化为 $z=z_0+R\ee^{\ii\theta}$, 则 $\dd z=\ii R\ee^{\ii\theta}\dd\theta$. 代入 $f(z_0)=(2\pi\ii)^{-1}\int_{C_R}f(z)/(z-z_0)\dd z$, 分子中的切向因子与分母抵消, 正好得到题给均值式.
\end{solution}
